% GENERATED FILE. Edit canonical lessons, then run npm run generate:books.
% canonical-source-hash: fnv1a64:7f9e6f08f3865386
% canonical-lessons: TE-C165-cindu, TE-C165-gorugu, TE-C165-bigincu, TE-C165-guccu, TE-C165-allu, TE-R165-first-pass-words-from-chapters-157-160, TE-R165-second-pass-words-from-chapters-161-165

\chapter{To spill, To shave, To tighten, To prick, To knit}
\label{ch:te-to-spill-to-shave-to-tighten-to-prick-to-knit}
\hlchaptermodality{165}

\begin{chapteropening}
\textbf{By the end of this chapter:} \emph{I can say to spill, to shave, to tighten, to prick and to knit.}

Complete the last lesson of chapter 165: I can say to spill, to shave, to tighten, to prick and to knit.
\end{chapteropening}

\section[\texorpdfstring{\te{చిందు} (cindu)}{చిందు (cindu)}]{\te{చిందు} (cindu) — to spill}

\label{lesson:TE-C165-cindu}



\begin{quote}
Before the new one: say the Telugu for to tear, then the Telugu for to wrap.
\end{quote}

\subsection*{You'll want to know: \te{చిందు}}
\textbf{\te{చిందు}} — \emph{cindu} — \textquotedblleft{}to spill\textquotedblright{}.

\begin{grammarlens}[title={What you've built}]
One more word for this chapter.
\end{grammarlens}

\subsection*{Your turn}
\begin{itemize}
\raggedright
  \item \emph{Say it:} \emph{cindu}
  \item \emph{Say it:} \emph{cindu}, once more
  \item \emph{Say it:} say \emph{cindu}
  \item \emph{Recall:} say the Telugu for to tear, then the Telugu for to wrap, then say \emph{cindu} again
\end{itemize}

\begin{tcolorbox}[breakable,colback=teal!4,colframe=teal!35!black,title={Before you move on}]
What does \te{చిందు} mean? (\textquotedblleft{}To spill\textquotedblright{}.) Say it once more.
\end{tcolorbox}
\section[\texorpdfstring{\te{గొరుగు} (gorugu)}{గొరుగు (gorugu)}]{\te{గొరుగు} (gorugu) — to shave}

\label{lesson:TE-C165-gorugu}



\begin{quote}
Before the new one: say the Telugu for to wrap, then the Telugu for to spill.
\end{quote}

\subsection*{You'll want to know: \te{గొరుగు}}
\textbf{\te{గొరుగు}} — \emph{gorugu} — \textquotedblleft{}to shave\textquotedblright{}.

\begin{grammarlens}[title={What you've built}]
One more word for this chapter.
\end{grammarlens}

\subsection*{Your turn}
\begin{itemize}
\raggedright
  \item \emph{Say it:} \emph{gorugu}
  \item \emph{Say it:} \emph{gorugu}, once more
  \item \emph{Say it:} say \emph{gorugu}
  \item \emph{Recall:} say the Telugu for to wrap, then the Telugu for to spill, then say \emph{gorugu} again
\end{itemize}

\begin{tcolorbox}[breakable,colback=teal!4,colframe=teal!35!black,title={Before you move on}]
What does \te{గొరుగు} mean? (\textquotedblleft{}To shave\textquotedblright{}.) Say it once more.
\end{tcolorbox}
\section[\texorpdfstring{\te{బిగించు} (bigiñcu)}{బిగించు (bigiñcu)}]{\te{బిగించు} (bigiñcu) — to tighten}

\label{lesson:TE-C165-bigincu}



\begin{quote}
Before the new one: say the Telugu for to spill, then the Telugu for to shave.
\end{quote}

\subsection*{You'll want to know: \te{బిగించు}}
\textbf{\te{బిగించు}} — \emph{bigiñcu} — \textquotedblleft{}to tighten\textquotedblright{}.

\begin{grammarlens}[title={What you've built}]
One more word for this chapter.
\end{grammarlens}

\subsection*{Your turn}
\begin{itemize}
\raggedright
  \item \emph{Say it:} \emph{bigiñcu}
  \item \emph{Say it:} \emph{bigiñcu}, once more
  \item \emph{Say it:} say \emph{bigiñcu}
  \item \emph{Recall:} say the Telugu for to spill, then the Telugu for to shave, then say \emph{bigiñcu} again
\end{itemize}

\begin{tcolorbox}[breakable,colback=teal!4,colframe=teal!35!black,title={Before you move on}]
What does \te{బిగించు} mean? (\textquotedblleft{}To tighten\textquotedblright{}.) Say it once more.
\end{tcolorbox}
\section[\texorpdfstring{\te{గుచ్చు} (guccu)}{గుచ్చు (guccu)}]{\te{గుచ్చు} (guccu) — to prick, to string (flowers)}

\label{lesson:TE-C165-guccu}



\begin{quote}
Before the new one: say the Telugu for to shave, then the Telugu for to tighten.
\end{quote}

\subsection*{You'll want to know: \te{గుచ్చు}}
\textbf{\te{గుచ్చు}} — \emph{guccu} — \textquotedblleft{}to prick, to string (flowers)\textquotedblright{}.

\begin{grammarlens}[title={What you've built}]
One more word for this chapter.
\end{grammarlens}

\subsection*{Your turn}
\begin{itemize}
\raggedright
  \item \emph{Say it:} \emph{guccu}
  \item \emph{Say it:} \emph{guccu}, once more
  \item \emph{Say it:} say \emph{guccu}
  \item \emph{Recall:} say the Telugu for to shave, then the Telugu for to tighten, then say \emph{guccu} again
\end{itemize}

\begin{tcolorbox}[breakable,colback=teal!4,colframe=teal!35!black,title={Before you move on}]
What does \te{గుచ్చు} mean? (\textquotedblleft{}To prick, to string (flowers)\textquotedblright{}.) Say it once more.
\end{tcolorbox}
\section[\texorpdfstring{\te{అల్లు} (allu)}{అల్లు (allu)}]{\te{అల్లు} (allu) — to knit, to weave}

\label{lesson:TE-C165-allu}



\begin{quote}
Before the new one: say the Telugu for to tighten, then the Telugu for to prick.
\end{quote}

\subsection*{You'll want to know: \te{అల్లు}}
\textbf{\te{అల్లు}} — \emph{allu} — \textquotedblleft{}to knit, to weave\textquotedblright{}.

\begin{grammarlens}[title={What you've built}]
One more word for this chapter.
\end{grammarlens}

\subsection*{Your turn}
\begin{itemize}
\raggedright
  \item \emph{Say it:} \emph{allu}
  \item \emph{Say it:} \emph{allu}, once more
  \item \emph{Say it:} say \emph{allu}
  \item \emph{Recall:} say the Telugu for to tighten, then the Telugu for to prick, then say \emph{allu} again
\end{itemize}

\begin{tcolorbox}[breakable,colback=teal!4,colframe=teal!35!black,title={Before you move on}]
What does \te{అల్లు} mean? (\textquotedblleft{}To knit, to weave\textquotedblright{}.) Say it once more.
\end{tcolorbox}
\section[(dialogue)]{Words from chapters 157-160 — a first pass}

\label{lesson:TE-R165-first-pass-words-from-chapters-157-160}



\begin{quote}
Say the words for to prick and to knit.
\end{quote}

\subsection*{Your turn}
Say these aloud:

\begin{itemize}
\raggedright
  \item to pay, to fry, to boil, to mix, to grind
  \item to hurry, to brush, to change, to take off, to measure
  \item to squeeze, to soak, to repair, to throw away, to lose
  \item to earn, to explain, to translate, to remember, to agree
\end{itemize}

\begin{tcolorbox}[breakable,colback=teal!4,colframe=teal!35!black,title={Before you move on}]
To pay? (\textbf{\te{చెల్లించు}}, \emph{celliñcu}.) To fry? (\textbf{\te{వేయించు}}, \emph{vēyiñcu}.) To boil? (\textbf{\te{ఉడికించు}}, \emph{uḍikiñcu}.) To mix? (\textbf{\te{కలుపు}}, \emph{kalupu}.) To grind? (\textbf{\te{రుబ్బు}}, \emph{rubbu}.) To hurry? (\textbf{\te{తొందరపడు}}, \emph{tondarapaḍu}.) To brush? (\textbf{\te{తోము}}, \emph{tōmu}.) To change? (\textbf{\te{మార్చు}}, \emph{mārcu}.) To take off? (\textbf{\te{విప్పు}}, \emph{vippu}.) To measure? (\textbf{\te{కొలుచు}}, \emph{kolucu}.) To squeeze? (\textbf{\te{పిండు}}, \emph{piṇḍu}.) To soak? (\textbf{\te{నానబెట్టు}}, \emph{nānabeṭṭu}.) To repair? (\textbf{\te{బాగుచేయు}}, \emph{bāgucēyu}.) To throw away? (\textbf{\te{పారేయు}}, \emph{pārēyu}.) To lose? (\textbf{\te{పోగొట్టుకో}}, \emph{pōgoṭṭukō}.) To earn? (\textbf{\te{సంపాదించు}}, \emph{sampādiñcu}.) To explain? (\textbf{\te{వివరించు}}, \emph{vivariñcu}.) To translate? (\textbf{\te{అనువదించు}}, \emph{anuvadiñcu}.) To remember? (\textbf{\te{గుర్తుపెట్టుకో}}, \emph{gurtupeṭṭukō}.) To agree? (\textbf{\te{ఒప్పుకో}}, \emph{oppukō}.)
\end{tcolorbox}
\section[(dialogue)]{Words from chapters 161-165 — a second pass}

\label{lesson:TE-R165-second-pass-words-from-chapters-161-165}



\begin{quote}
Say the words for to dance and to celebrate.
\end{quote}

\subsection*{Your turn}
Say these aloud:

\begin{itemize}
\raggedright
  \item to dance, to celebrate, to greet, to feel shy, to be surprised
  \item to kiss, to hug, to tease, to praise, to describe
  \item to suspect, to realize, to pay attention to, to compete, to kick
  \item to grate, to strain, to dip, to tear, to wrap
  \item to spill, to shave, to tighten, to prick, to knit
\end{itemize}

\begin{tcolorbox}[breakable,colback=teal!4,colframe=teal!35!black,title={Before you move on}]
To dance? (\textbf{\te{నాట్యం} \te{చేయు}}, \emph{nāṭyaṁ cēyu}.) To celebrate? (\textbf{\te{జరుపుకో}}, \emph{jarupukō}.) To greet? (\textbf{\te{పలకరించు}}, \emph{palakariñcu}.) To feel shy? (\textbf{\te{సిగ్గుపడు}}, \emph{siggupaḍu}.) To be surprised? (\textbf{\te{ఆశ్చర్యపోవు}}, \emph{āścaryapōvu}.) To kiss? (\textbf{\te{ముద్దుపెట్టు}}, \emph{muddupeṭṭu}.) To hug? (\textbf{\te{కౌగిలించు}}, \emph{kaugiliñcu}.) To tease? (\textbf{\te{ఆటపట్టించు}}, \emph{āṭapaṭṭiñcu}.) To praise? (\textbf{\te{పొగడు}}, \emph{pogaḍu}.) To describe? (\textbf{\te{వర్ణించు}}, \emph{varṇiñcu}.) To suspect? (\textbf{\te{అనుమానించు}}, \emph{anumāniñcu}.) To realize? (\textbf{\te{గ్రహించు}}, \emph{grahiñcu}.) To pay attention to? (\textbf{\te{పట్టించుకో}}, \emph{paṭṭiñcukō}.) To compete? (\textbf{\te{పోటీపడు}}, \emph{pōṭīpaḍu}.) To kick? (\textbf{\te{తన్ను}}, \emph{tannu}.) To grate? (\textbf{\te{తురుము}}, \emph{turumu}.) To strain? (\textbf{\te{వడకట్టు}}, \emph{vaḍakaṭṭu}.) To dip? (\textbf{\te{ముంచు}}, \emph{muñcu}.) To tear? (\textbf{\te{చింపు}}, \emph{cimpu}.) To wrap? (\textbf{\te{చుట్టు}}, \emph{cuṭṭu}.) To spill? (\textbf{\te{చిందు}}, \emph{cindu}.) To shave? (\textbf{\te{గొరుగు}}, \emph{gorugu}.) To tighten? (\textbf{\te{బిగించు}}, \emph{bigiñcu}.) To prick? (\textbf{\te{గుచ్చు}}, \emph{guccu}.) To knit? (\textbf{\te{అల్లు}}, \emph{allu}.)
\end{tcolorbox}
